\subsection{Ample comparison quotients on Hirzebruch surfaces}
\label{subsec:hirzebruch-ample-comparison}
\index{Hirzebruch surface!ample comparison}

To apply \cref{cor:gauduchon-semipositive-finsler-obstruction}, we need ample vector bundles whose duals are close to a split bundle $F_1\oplus F_2$ with $\deg_{\Theta_s}F_s>0$ for suitable test forms $\Theta_s$. The split bundle $F_1^*\oplus F_2^*$ is not ample itself: otherwise its quotient line bundles $F_s^*$ would be ample, hence carry positively curved metrics, and each $F_s$ would have negative degree for every test form. The ample bundles must therefore be obtained by deformation. On a Hirzebruch surface we choose a split bundle $E_0=L_1\oplus L_2$ whose dual summands have positive degree for suitable K\"ahler forms, and we construct an ample bundle $E_{\mathrm{amp}}$ with the same Chern classes as a quotient of a trivial bundle. Both bundles belong to the family of bundles with these Chern classes that are prioritary for the ruling. This family is irreducible, which will allow us to find ample bundles in every neighborhood of $E_0$.

Let $e\geq0$ be an integer, and let
\[
 X=\mathbb F_e=\mathbb P_{\mathbb P^1}(\mathcal O\oplus\mathcal O(e)),
 \qquad \pi\colon X\longrightarrow\mathbb P^1,
\]
be the Hirzebruch surface, where the projective bundle parametrizes lines. Write $S$ for the section defined by the subbundle $\mathcal O(e)$ and $f$ for a fiber of $\pi$, called a ruling fiber; we identify line bundles with divisor classes and write $\mathcal O_X(aS+bf)$. The normal bundle of $S$ is $\mathcal O(-e)$, so
\[
 S^2=-e,\qquad S\cdot f=1,\qquad f^2=0.
\]
The dual of the tautological line bundle is $\mathcal O_X(S)$. A line bundle of degree zero on the ruling fibers is pulled back from $\mathbb P^1$: its direct image has rank one, and the evaluation map from the pullback of the direct image is an isomorphism on each fiber. Consequently $\operatorname{Pic}(X)=\mathbb ZS\oplus\mathbb Zf$.

\Needspace{5\baselineskip}
\begin{lemma}\label{lem:hirzebruch-ruling-sections}
\index{Hirzebruch surface!ruling cohomology}
For integers $a\geq0$ and $b$,
\[
 \pi_*\mathcal O_X(aS+bf)=
 \bigoplus_{i=0}^a\mathcal O_{\mathbb P^1}(b-ie),
 \qquad R^1\pi_*\mathcal O_X(aS+bf)=0.
\]
A class $aS+bf$ is effective exactly when $a,b\geq0$. If $a\geq1$ and $b>ae$, its line bundle is very ample. In particular,
\[
 H=S+(3e+1)f,
 \qquad H^2=5e+2,
 \qquad h^0(X,\mathcal O_X(H))=5e+4.
\]
The canonical class is $K_X=-2S-(e+2)f$.
\end{lemma}
\begin{proof}
Over a trivializing chart of the base, sections of $\mathcal O_X(aS)$ are homogeneous polynomials of degree $a$ in the fiber coordinates. Their transition law gives
\[
 \pi_*\mathcal O_X(aS)=
 \operatorname{Sym}^a(\mathcal O\oplus\mathcal O(-e)),
\]
and twisting by the base gives the formula. The first cohomology of $\mathcal O_{\mathbb P^1}(a)$ on each fiber vanishes, which gives the assertion on the higher direct image. If $a<0$, restriction to a general ruling fiber excludes a nonzero section. For $a\geq0$, the direct-image formula gives a nonzero section exactly when $b\geq0$.

Suppose that $a\geq1$ and $b>ae$. All line bundles in the direct image then have degree at least one. Fiber polynomials of degree $a$ separate points and tangent directions within a fiber. To separate points over different base points, choose a fiber monomial that is nonzero at the first point and a coefficient section vanishing at the second base point but not at the first. At any point, choose such a nonzero monomial and a coefficient section with a simple zero at its base point; the resulting global section vanishes there, has nonzero derivative in a prescribed base direction and has zero vertical derivative. The complete linear system therefore separates points and first jets, so it defines an embedding. The formulas for $H^2$ and $h^0(H)$ follow from the intersection numbers and from the direct image $\mathcal O(3e+1)\oplus\mathcal O(2e+1)$. Finally, the relative Euler sequence gives $K_{X/\mathbb P^1}=-2S-ef$, and adding the pullback of the canonical class of $\mathbb P^1$ proves the last identity.
\end{proof}

Put
\[
 L_1=\mathcal O_X(5S-f),\qquad
 L_2=\mathcal O_X(-S+(12e+5)f),\qquad
 E_0=L_1\oplus L_2.
\]
Their Chern classes satisfy
\[
 c_1(E_0)=4H,\qquad c_2(E_0)=65e+26=13H^2.
\]
An ample bundle with these Chern classes can be constructed as a quotient of a trivial bundle. The main point of the proof is to make the fibers of the associated map to projective space finite, which requires excluding common curves of the defining sections uniformly over the projective parameter.

\Needspace{5\baselineskip}
\begin{proposition}\label{prop:hirzebruch-ample-comparison}
\index{Vector bundle!ample comparison quotient}
There exist linearly independent sections $a_0,\ldots,a_3$ of $\mathcal O_X(H)$ with no common zero. For any such quadruple and a general quadruple $b_0,\ldots,b_3$ of sections of $\mathcal O_X(3H)$, the sequence
\[
 0\longrightarrow\mathcal O_X(-H)\oplus\mathcal O_X(-3H)
 \xrightarrow{(a,b)}\mathcal O_X^{\oplus4}
 \longrightarrow E_{\mathrm{amp}}\longrightarrow0
\]
has a locally free ample cokernel $E_{\mathrm{amp}}$ of rank two, with $c_1(E_{\mathrm{amp}})=4H$ and $c_2(E_{\mathrm{amp}})=13H^2$. Here \emph{general} means outside a proper Zariski-closed subset of the space of quadruples.
\end{proposition}
\begin{proof}
\emph{Step 1: the quadruple $a$.} Since $\mathcal O_X(H)$ is globally generated, a common zero at a fixed point imposes four independent linear conditions on a quadruple $a$. The incidence variety of pairs consisting of a point of the surface and a quadruple vanishing there therefore has dimension at most two less than that of the space of quadruples, and its image in that space is closed by properness. Hence quadruples with a common zero form a proper closed subset. Linear independence is also a nonempty open condition, because $h^0(H)\geq4$, and the two open conditions intersect. Fix such an $a$ and write
\[
 V=H^0(X,\mathcal O_X(3H))^{\oplus4},\qquad N=\dim V.
\]

\emph{Step 2: local freeness.} At a fixed point $x$, the map $(a,b)$ drops rank exactly when $b(x)$ is proportional to the nonzero vector $a(x)$. Since evaluation on $V$ is surjective, this is a condition of codimension three. Its incidence over $X$ has dimension at most $N-1$, and its image in $V$ is closed. Outside that image the cokernel is locally free of rank two.

\emph{Step 3: the map to projective space.} Dualizing the sequence identifies $E_{\mathrm{amp}}^*$ with a subbundle of $(\mathbb C^4)^*\otimes\mathcal O_X$. This inclusion defines a proper morphism
\[
 q_b\colon\mathbb P(E_{\mathrm{amp}}^*)\longrightarrow
 \mathbb P((\mathbb C^4)^*),
 \qquad \mathcal O_{\mathbb P(E_{\mathrm{amp}}^*)}(1)
 =q_b^*\mathcal O_{\mathbb P^3}(1).
\]
For $w=[w_0:\cdots:w_3]$, its fiber is the common zero scheme of
\[
 s_w=\sum_iw_i a_i\in H^0(X,\mathcal O_X(H)),
 \qquad t_w=\sum_iw_i b_i\in H^0(X,\mathcal O_X(3H)).
\]
Linear independence of $a$ gives $s_w\neq0$ for every $w$.

\emph{Step 4: no common curves.} A nonzero effective subdivisor $D$ of $\operatorname{div}(s_w)$ has class $uS+vf$, where
\[
 u\in\{0,1\},\qquad 0\leq v\leq3e+1,
 \qquad (u,v)\neq(0,0).
\]
Write $|D|=\mathbb P H^0(X,\mathcal O_X(D))$ for the complete linear system. For each of these finitely many classes $D_0$, let $T_{D_0}$ parametrize the triples $(w,[\sigma],[\tau])$ with
\[
 [\sigma]\in|D_0|,\qquad [\tau]\in|H-D_0|,
 \qquad [s_w]=[\sigma\tau].
\]
Since products of nonzero sections are nonzero, this is a closed projective parameter scheme. Its fiber over $w$ consists of the finitely many effective subdivisors of the fixed divisor $\operatorname{div}(s_w)$, counted with their possible multiplicities. Thus $\dim T_{D_0}\leq3$.

Vanishing on a nonzero effective divisor $D$ imposes at least four independent conditions on $H^0(X,\mathcal O_X(3H))$. Indeed, choose an integral curve $C$ in its support. Very ampleness gives sections $r_0,r_1$ of $\mathcal O_X(H)$ with nonconstant ratio on $C$. The four cubic monomials in their restrictions are linearly independent, since a linear relation would make the nonconstant ratio algebraic over $\mathbb C$, hence constant.

For fixed $(w,[\sigma],[\tau])$, the map $V\to H^0(X,\mathcal O_X(3H))$, $b\mapsto t_w$, is surjective. The condition that $t_w$ vanish on $D=\operatorname{div}(\sigma)$ therefore has codimension at least four in $V$. This condition is closed in families: the universal multiplication map
\[
 H^0(X,\mathcal O_X(3H-D_0))\otimes\mathcal O_{|D_0|}(-1)
 \longrightarrow H^0(X,\mathcal O_X(3H))\otimes\mathcal O_{|D_0|}
\]
is injective on fibers, so its image is a subbundle, and the condition is the vanishing of the image of $t_w$ in the quotient bundle. Its incidence over $T_{D_0}$ has dimension at most $N-1$, and its image in $V$ is closed.

\emph{Step 5: conclusion.} Choose $b$ outside these finitely many proper closed images and outside the image of Step~2. Then $s_w$ and $t_w$ have no common curve for any $w$, so every fiber of $q_b$ is finite. A proper morphism with finite fibers is finite, and the pullback of $\mathcal O_{\mathbb P^3}(1)$ by a finite morphism is ample; hence $E_{\mathrm{amp}}$ is ample. Finally, the Whitney formula gives
\[
 c(E_{\mathrm{amp}})=\frac1{(1-H)(1-3H)}=1+4H+13H^2.
\]
\end{proof}

\index{Prioritary bundle}
A vector bundle $E$ on $X$ is called \emph{prioritary for the ruling} if $\operatorname{Ext}^2_X(E,E(-f))=0$. Since all ruling fibers are linearly equivalent, this condition does not depend on the chosen fiber, and by Serre duality it is equivalent to
\[
 \operatorname{Hom}_X(E,E(K_X+f))=0.
\]
Prioritary bundles of fixed rank and Chern classes form an irreducible family, by the theorem of Walter recalled in \cref{subsec:hirzebruch-stable-ample-deformations}.

\Needspace{5\baselineskip}
\begin{lemma}\label{lem:hirzebruch-prioritary-comparison}
\index{Prioritary bundle!comparison points}
Both $E_0$ and $E_{\mathrm{amp}}$ are prioritary for the ruling.
\end{lemma}
\begin{proof}
Write $D=K_X+f=-2S-(e+1)f$. The summands of $\operatorname{End}(E_0)(D)$ have classes
\[
 D,\quad D,\quad 4S-(13e+7)f,\quad -8S+(11e+5)f.
\]
Each has a negative coefficient in the basis $S,f$, so by \cref{lem:hirzebruch-ruling-sections} none of them has a nonzero section. This proves the required vanishing of $\operatorname{Hom}$ for $E_0$.

Precomposition with the surjection $\mathcal O_X^4\to E_{\mathrm{amp}}$ gives an injection
\[
 \operatorname{Hom}(E_{\mathrm{amp}},E_{\mathrm{amp}}(D))
 \hookrightarrow H^0(X,E_{\mathrm{amp}}(D))^4.
\]
Twist the defining sequence by $D$. Its middle term has no sections. For $j=1,3$, Serre duality and the formula of \cref{lem:hirzebruch-ruling-sections} give
\[
 \begin{aligned}
 H^1(X,\mathcal O_X(D-jH))&=H^1(X,\mathcal O_X(jH-f))^*,\\
 \pi_*\mathcal O_X(jH-f)&=\bigoplus_{i=0}^j\mathcal O_{\mathbb P^1}(j(3e+1)-1-ie).
 \end{aligned}
\]
The higher direct image vanishes, and every degree in this sum is at least $j(2e+1)-1\geq0$. Hence all these first cohomology groups vanish, and the long exact sequence of the twisted sequence gives $H^0(X,E_{\mathrm{amp}}(D))=0$. This proves the vanishing of $\operatorname{Hom}$ for $E_{\mathrm{amp}}$.
\end{proof}

\Needspace{5\baselineskip}
\begin{proposition}\label{prop:hirzebruch-positive-dual-tests}
\index{Gauduchon degree!Hirzebruch dual tests}
For $F_s=L_s^*$, $s=1,2$, there are K\"ahler forms $\eta_s$ with $\deg_{\eta_s}F_s=1$. Consequently, the split bundle $E_0^*=F_1\oplus F_2$ satisfies the hypotheses of \cref{thm:convex-norm-deformation-obstruction}.
\end{proposition}
\begin{proof}
The integral classes
\[
 A_1=6S+(6e+1)f,\qquad A_2=S+(13e+6)f
\]
are very ample by \cref{lem:hirzebruch-ruling-sections}. Pullbacks of normalized Fubini--Study forms under the corresponding embeddings give K\"ahler forms $\eta_s$ representing $c_1(\mathcal O_X(A_s))$. A direct computation of intersection numbers gives
\[
 c_1(L_1)\cdot A_1=-1,\qquad
 c_1(L_2)\cdot A_2=-1,
\]
so the degrees of the duals are equal to one. Since $n=2$, the closed positive $(1,1)$-forms $\eta_s$ are admissible test forms $\Theta_s$ in \cref{thm:convex-norm-deformation-obstruction}.
\end{proof}
